% Date : 02/11/2022
% Version : 1
% Auteur : Erwan Capitaine
% Relu : 
% Relecteur : 

% ------------------------------------------------------------ %
% ----------------------  Fiche MAG01  ----------------------- %
% Thème : Électromagnétisme - champ électrique
% Classes : Toutes les classes de sup
% ------------------------------------------------------------ %



% ======================================================= % 
% ============== Gestion de la compilation ============== %
% ======================================================= % 
\ifdefined\mainIsLoaded\else
\RequirePackage{import}
\subimport{../../}{_preambule_CdE_PC}
\begin{document}
\fi
% ======================================================= % 
% ======================================================= % 


% ======================================================= % 
% ===================== Couleurs ======================== %
% ======================================================= % 

% -----------------------------------%
%  Définition de la couleur A de CBD
% Merci de préfixer vos couleur

% La commande \couleurNBouCouleur prend deux paramètres :
% #1 est la couleur NB ; #2 est la couleur "en couleur"

\def\CBDcouleurA{\couleurNBouCouleur{black}{\myCBDBlue}}

% -----------------------------------%

% ======================================================= % 
% ======================================================= % 



% ============================================================================ %
%                            FICHE D'ENTRAÎNEMENT                              %
% ============================================================================ %
% ======================= Métadonnées sur la fiche =========================== %
\titreFicheEntrainement		{Champ électrique}
\grandTheme					{Électromagnétisme}
\numeroFiche				{MAG01}
\uniqueID					{tDpmMaBgLyL} % une chaîne de caractères (a-z, A-Z) aléatoire de 10 caractères.
\nombreColonnesReponses		{2} % 1, 2, 3, etc.
% ============================================================================ %

%%%%%%%%%%%%%%%%%%%%%%%%%%%%%%%%%%%%%%%%%%%%%%%%%%%%%%%%%%%%%%%%%%%%%%%%%%%%%%%%%%%%%%%%%%%%%%
\debutFicheEntrainement                                                                      %
%%%%%%%%%%%%%%%%%%%%%%%%%%%%%%%%%%%%%%%%%%%%%%%%%%%%%%%%%%%%%%%%%%%%%%%%%%%%%%%%%%%%%%%%%%%%%%




% --------------------------------------------------------------------- %
%                              Prérequis                                %
%                             (facultatif)                              %
% --------------------------------------------------------------------- %
% ================== Métadonnées sur le prérequis ===================== %
\avecPrerequis{Y} % Y ou N
% ===================================================================== %
\begin{prerequis}
	Projections. Coordonnées polaires. Développement limité. \\
	Dérivation et intégration.
\end{prerequis}
% --------------------------------------------------------------------- %






% •••••••••••••••••••••••••••••••••••••••••••••••••••••••••••••••••••••••••••• %
% •••••••••••••••••••••••••••••••••••••••••••••••••••••••••••••••••••••••••••• %
\sectionFicheEntrainement{Pour commencer}
% •••••••••••••••••••••••••••••••••••••••••••••••••••••••••••••••••••••••••••• %
% •••••••••••••••••••••••••••••••••••••••••••••••••••••••••••••••••••••••••••• %



% ********************************************************************* %
%                           ENTRAÎNEMENT                                % 
% ********************************************************************* %
% ================ Métadonnées sur l'entraînement ===================== %

\titreEntrainementFacultatif	{Projection d'une force}
\hauteurLargeurCadreReponse		{8mm}{4cm}
\dureeResolutionFacultative		{2} % 1, 2, 3 ou 4
\typeEntrainement				{Newton} % Newton ou QCM ou AN ou vide (exercice "normal")
\basiqueEtTransversal			{Y}
\nombreColonnesQuestions		{1} % vide, 1, 2, 3, etc.
\avecPlusieursQuestions			{Y} % Y ou N
\initialisationEntrainement
% ===================================================================== %


Une charge électrique $q$ située en un point $\ptB\left(a,0\right)$ exerce une force $\vv{F}$ sur une autre charge $q_0$ située au point $\ptA\left( 0,y \right)$.
\begin{center}
	\subimport{_images/}{fiche_MAG01_figure_projection_force_CBD_v3.tex}
\end{center}


% =============================================================================== %
\debutEntrainement
% =============================================================================== %



% ^^^^^^^^^^^^^^^^^^^^^^^^^^^^^^^^^^^^^^ %
%               Question                 %
% ^^^^^^^^^^^^^^^^^^^^^^^^^^^^^^^^^^^^^^ %

\begin{enonce}
Exprimer la distance $\ptB\ptA$ en fonction de $a$ et de $y$. \minusBigskip

\end{enonce}

\reponse{$\sqrt{a^2 + y^2}$}

\begin{corrige}
Dans le triangle rectangle $\ptO\ptA\ptB$, on a ${\ptB\ptA} = \sqrt{a^2 + y^2}$.
\end{corrige}

% ************************************** %


% ^^^^^^^^^^^^^^^^^^^^^^^^^^^^^^^^^^^^^^ %
%               Question                 %
% ^^^^^^^^^^^^^^^^^^^^^^^^^^^^^^^^^^^^^^ %

\begin{enonce}
Exprimer $\cos (\alpha)$ en fonction de $a$ et $y$. \minusBigskip

\end{enonce}

\reponse{$\frac{a}{\sqrt{a^2 + y^2}}$}

\begin{corrige}
Dans le triangle rectangle $\ptO\ptA\ptB$, on a $\cos (\alpha) = \frac{a}{\ptB\ptA} = \frac{a}{\sqrt{a^2 + y^2}}$.
\end{corrige}

% ************************************** %



% ^^^^^^^^^^^^^^^^^^^^^^^^^^^^^^^^^^^^^^ %
%               Question                 %
% ^^^^^^^^^^^^^^^^^^^^^^^^^^^^^^^^^^^^^^ %

\begin{enonce}
Exprimer $\sin (\alpha)$ en fonction de $a$ et $y$.\minusBigskip

\end{enonce}

\reponse{$\frac{y}{\sqrt{a^2 + y^2}}$ }

\begin{corrige}
Dans le triangle rectangle $\ptO\ptA\ptB$, on a 
	$\sin (\alpha) = \frac{y}{\ptB\ptA} = \frac{y}{\sqrt{a^2 + y^2}}$.
\end{corrige}

% ************************************** %




% ^^^^^^^^^^^^^^^^^^^^^^^^^^^^^^^^^^^^^^ %
%               Question                 %
% ^^^^^^^^^^^^^^^^^^^^^^^^^^^^^^^^^^^^^^ %

\begin{enonce}
Décomposer le vecteur $\vv{F}$ dans la base $(\vv{e_x}, \vv{e_y})$ %et exprimer ses composantes 
en fonction de sa norme $\lVert \vv{F} \rVert$, $a$ et $y$.

\end{enonce}

\reponse{$\frac{\lVert \vv{F} \rVert}{\sqrt{a^2 + y^2}}\left( -a \vv{e_x} + y \vv{e_y} \right)$}

\begin{corrige}
La composante suivant $\vv{e_x}$ correspond au produit scalaire
\[
	{F}_x=\vv{F}\cdot \vv{e_x}=\lVert \vv{F}\rVert\cos(\alpha+\pi)=-\lVert \vv{F} \rVert\cos (\alpha).
\]
De même, la composante suivant $\vv{e_y}$ correspond à
\[
	{F}_y=\vv{F}\cdot \vv{e_y}=\lVert \vv{F} \rVert\cos\left(-\frac{\pi}{2}+\alpha\right)=\lVert \vv{F} \rVert\sin\alpha.
	\]
Ainsi, on a
\[
{F}_x=-\lVert \vv{F}\rVert\frac{a}{\sqrt{a^2 + y^2}} \quad \text{et} \quad {F}_y=\lVert \vv{F} \rVert\frac{y}{\sqrt{a^2 + y^2}}.
\]
Finalement, on a 
\[
\vv{F}= {F}_x \vv{e_x} + {F}_y \vv{e_y} = \frac{\lVert \vv{F} \rVert}{\sqrt{a^2 + y^2}}\left( -a \vv{e_x} + y \vv{e_y} \right).
\]
\end{corrige}

% ************************************** %

% =============================================================================== %
\finEntrainement
% =============================================================================== %








% ********************************************************************* %
%                           ENTRAÎNEMENT                                % 
% ********************************************************************* %
% ================ Métadonnées sur l'entraînement ===================== %

\titreEntrainementFacultatif	{Un combat d'interaction électrique}
\hauteurLargeurCadreReponse		{6mm}{6cm}
\dureeResolutionFacultative		{1} % 1, 2, 3 ou 4
\typeEntrainement				{QCM} % Newton ou QCM ou AN ou vide (exercice "normal")
\basiqueEtTransversal			{Y}
\nombreColonnesQuestions		{1} % vide, 1, 2, 3, etc.
\avecPlusieursQuestions			{N} % Y ou N
\initialisationEntrainement
% ===================================================================== %

% =============================================================================== %
\debutEntrainement
% =============================================================================== %

% ^^^^^^^^^^^^^^^^^^^^^^^^^^^^^^^^^^^^^^ %
%               Question                 %
% ^^^^^^^^^^^^^^^^^^^^^^^^^^^^^^^^^^^^^^ %

\begin{enonce}
	On étudie une charge électrique $q_0$ positive. La valeur de la force $F$ qu'exerce une autre charge $q$ sur $q_0$ est telle que $F = C \frac{q}{d^2}$ où $d$ est la distance entre les deux charges et où $C$ est une constante.

	Laquelle de ces quatre charges attire le plus fortement la charge $q_0$ ?
	
	\begin{listeQCM2Colonnes}
	\item $q = \np[C]{2,00}$ et $ d = \np[mm]{4,00}$
	\item $q = \np[kC]{-5,0}$ et $ d = \np[m]{0,4}$
	\item $q = \np[mC]{-3,0}$ et $ d = \np[\mu m]{200}$
	\item $q = \np[C]{100}$ et $ d = \np[cm]{20}$
	\end{listeQCM2Colonnes}

	\bigskip

	\vphantom{Calculer}
\end{enonce}

\reponse{ \reponseC{}}

\begin{corrige}
	Une force attractive a une valeur négative, la charge qui attire le plus est donc la charge avec la force négative la plus importante en valeur absolue, soit la réponse \reponseC{}. En effet, on a

	\begin{listeQCM2Colonnes}
		\item $F/C = \frac{ \np[C]{2,00} }{\left(\np[m]{4,00e{-3}}\right)^2} = \np[C.m^{-2}]{1,25e{5}}$
		\item $F/C = \frac{ \np[C]{-5,0e{3}} }{\left(\np[m]{0,4}\right)^2} = \np[C.m^{-2}]{-3,1e{4}}$
		\item $F/C = \frac{ \np[C]{-3,0e{-3}} }{\left(\np[m]{200e{-6}}\right)^2} = \np[C.m^{-2}]{-7,5e{4}}$
		\item $F/C = \frac{ \np[C]{100} }{\left(\np[m]{20e{-2}}\right)^2} = \np[C.m^{-2}]{2,5e{3}}$
	\end{listeQCM2Colonnes}

\end{corrige}

% ************************************** %

% =============================================================================== %
\finEntrainement
% =============================================================================== %






\newpage




% •••••••••••••••••••••••••••••••••••••••••••••••••••••••••••••••••••••••••••• %
% •••••••••••••••••••••••••••••••••••••••••••••••••••••••••••••••••••••••••••• %
\sectionFicheEntrainement{Étude de charges ponctuelles}
% •••••••••••••••••••••••••••••••••••••••••••••••••••••••••••••••••••••••••••• %
% •••••••••••••••••••••••••••••••••••••••••••••••••••••••••••••••••••••••••••• %




% ********************************************************************* %
%                           ENTRAÎNEMENT                                % 
% ********************************************************************* %
% ================ Métadonnées sur l'entraînement ===================== %

\titreEntrainementFacultatif	{Force due à deux charges}
\hauteurLargeurCadreReponse		{6mm}{3cm}
\dureeResolutionFacultative		{1} % 1, 2, 3 ou 4
\typeEntrainement				{} % Newton ou QCM ou AN ou vide (exercice "normal")
\basiqueEtTransversal			{Y}
\nombreColonnesQuestions		{2} % vide, 1, 2, 3, etc.
\avecPlusieursQuestions			{Y} % Y ou N
\initialisationEntrainement
% ===================================================================== %

% mmmmmmmmmmmmmmmmmmmmmmmmmmmmmmmmmmmmmmmmmmmmmmmmmmmmmmmmm %
% 		  Insertion d'une image en regard d'un texte
% mmmmmmmmmmmmmmmmmmmmmmmmmmmmmmmmmmmmmmmmmmmmmmmmmmmmmmmmm %
% --------------------------------------------------------- %
\pourcentageDeLaPartieAGauche   {0.66}
% --------------------------------------------------------- %
%                                                           %
% -------------------- Partie à gauche -------------------- %
%                                                           %
                                \initialisationPartieGauche % 
%                                                           %
% 

La loi de Coulomb permet d'exprimer la force $\vv{F}_{1/0}$ exercée par une charge $q_1$ située en un point $\ptB$ sur une charge $q_0$ située en un point $\ptA$ :
\[
	\vv{F}_{1/0} = \frac{1}{4\pi\varepsilon_0}\frac{q_0 q_1}{{\ptB\ptA}^2}\vv{e}_{\ptB\ptA}
\]
avec $\varepsilon_0$ la permittivité du vide et $\vv{e}_{\ptB\ptA}$ le vecteur unitaire munissant le segment~$\ptB\ptA$.


On étudie les forces $\vv{F}_{1/0}$ et $\vv{F}_{2/0}$ exercées respectivement par les charges $q_1$ et $q_2$ sur la charge $q_0$.% avec des distances telles que $\ptB\ptA = \ptC\ptA$.

% --------------------------------------------------------- %
%                                                           %
% -------------------- Partie à droite -------------------- %
%                                                           %
\initialisationPartieDroite %
%                                                           %
%                                                           %
\begin{center}
	\subimport{_images/}{fiche_MAG01_figure_force_deux_charges_CBD_v2.tex}
\end{center}
% --------------------------------------------------------- %
                               \finalisationDuPartageDePage %					
% --------------------------------------------------------- %

Selon les différentes valeurs des charges $q_0$, $q_1$ et $q_2$, déterminer si la résultante des forces $\vv{F} = \vv{F}_{1/0} + \vv{F}_{2/0}$ est orientée selon $\vv{e_x}$, $-\vv{e_x}$, $\vv{e_y}$ ou $-\vv{e_y}$.

% =============================================================================== %
\debutEntrainement
% =============================================================================== %

% ^^^^^^^^^^^^^^^^^^^^^^^^^^^^^^^^^^^^^^ %
%               Question                 %
% ^^^^^^^^^^^^^^^^^^^^^^^^^^^^^^^^^^^^^^ %

\begin{enonce}
	$q_0 = q_1 = q_2$
\end{enonce}

\reponse{$\vv{e_y}$}

\begin{corrige}
	On a $q_0q_1 = q^2$ et $q_0q_2 = q^2$ donc $\vv{F}_{1/0} = -F_x \vv{e_x} + F_y \vv{e_y}$ et $\vv{F}_{2/0} = F_x \vv{e_x} + F_y \vv{e_y}$. Ainsi, la somme des deux forces est $\vv{F} = 2F_y \vv{e_y}$.
\end{corrige}

% ************************************** %


% ^^^^^^^^^^^^^^^^^^^^^^^^^^^^^^^^^^^^^^ %
%               Question                 %
% ^^^^^^^^^^^^^^^^^^^^^^^^^^^^^^^^^^^^^^ %

\begin{enonce}
	$-q_0 = -q_1 = q_2$
\end{enonce}

\reponse{$-\vv{e_x}$}

\begin{corrigeNewpage}
	On a $q_0q_1 = q^2$ et $q_0q_2 = -q^2$ donc $\vv{F}_{1/0} = -F_x \vv{e_x} + F_y \vv{e_y}$ et $\vv{F}_{2/0} = -F_x \vv{e_x} - F_y \vv{e_y}$. Ainsi, la somme des deux forces est $\vv{F} = -2F_x \vv{e_x}$.
\end{corrigeNewpage}

% ************************************** %


% ^^^^^^^^^^^^^^^^^^^^^^^^^^^^^^^^^^^^^^ %
%               Question                 %
% ^^^^^^^^^^^^^^^^^^^^^^^^^^^^^^^^^^^^^^ %

\begin{enonce}
	$q_0 = -q_1 = q_2$
\end{enonce}

\reponse{$\vv{e_x}$}

\begin{corrige}
	On a $q_0q_1 = -q^2$ et $q_0q_2 = q^2$ donc $\vv{F}_{1/0} = F_x \vv{e_x} - F_y \vv{e_y}$ et $\vv{F}_{2/0} = F_x \vv{e_x} + F_y \vv{e_y}$. Ainsi, la somme des deux forces est $\vv{F} = 2F_x \vv{e_x}$.
\end{corrige}

% ************************************** %


% ^^^^^^^^^^^^^^^^^^^^^^^^^^^^^^^^^^^^^^ %
%               Question                 %
% ^^^^^^^^^^^^^^^^^^^^^^^^^^^^^^^^^^^^^^ %

\begin{enonce}
	$-\frac{1}{2}q_0 = q_1 = q_2$
\end{enonce}

\reponse{$-\vv{e_y}$}

\begin{corrige}
	On a $q_0q_1 = -2q^2$ et $q_0q_2 = -2q^2$ donc $\vv{F}_{1/0} = F_x \vv{e_x} - F_y \vv{e_y}$ et $\vv{F}_{2/0} = -F_x \vv{e_x} - F_y \vv{e_y}$. Ainsi, la somme des deux forces est $\vv{F} = -2F_y \vv{e_y}$.
\end{corrige}

% ************************************** %

% =============================================================================== %
\finEntrainement
% =============================================================================== %







% ********************************************************************* %
%                           ENTRAÎNEMENT                                % 
% ********************************************************************* %
% ================ Métadonnées sur l'entraînement ===================== %

\titreEntrainementFacultatif	{Charge accélérée}
\hauteurLargeurCadreReponse		{8mm}{2.5cm}
\dureeResolutionFacultative		{2} % 1, 2, 3 ou 4
\typeEntrainement				{} % Newton ou QCM ou AN ou vide (exercice "normal")
\basiqueEtTransversal			{Y}
\nombreColonnesQuestions		{1} % vide, 1, 2, 3, etc.
\avecPlusieursQuestions			{Y} % Y ou N
\initialisationEntrainement
% ===================================================================== %

On considère une particule de charge $q$ et de masse $m$ se déplaçant le long d'un axe $(\ptO x)$ sous l'action d'un champ de potentiel électrique $V(x)$.

On dispose de trois expressions de $V(x)$ dont une seule est homogène
\[
	\text{ \reponseA{} } V(x) = V_0 \left( \frac{1}{x} - \frac{1}{a} \right) \qquad
	\text{ \reponseB{} } V(x) = V_0 \left( 1 - \left(\frac{x}{a}\right)^2\right) \qquad
	\text{ \reponseC{} } V(x) = V_0 \left( a^2 - x^2\right).
\]


La vitesse $v(x)$ de la particule et le potentiel $V(x)$ en un point $x$ sont liés par la relation
\begin{equation}
	\frac{1}{2}mv(x)^2 + qV(x) = \mathrm{C^{te}}. \label{relation_v_V}
\end{equation}

En $x=0$ la vitesse de la particule est nulle.

% =============================================================================== %
\debutEntrainement
% =============================================================================== %



% ^^^^^^^^^^^^^^^^^^^^^^^^^^^^^^^^^^^^^^ %
%               Question                 %
% ^^^^^^^^^^^^^^^^^^^^^^^^^^^^^^^^^^^^^^ %

\begin{enonce}
	Déterminer la seule expression de $V(x)$ homogène à un potentiel électrique
\end{enonce}

\reponse{\reponseB{}}

\begin{corrige}
	Comme $V_0$ est homogène à un potentiel électrique, l'argument entre parenthèses doit être sans dimension, ce qui est le cas dans l'expression
	\[
		V(x) = V_0 \left( 1 - \left(\frac{x}{a}\right)^2\right).	
	\]
\end{corrige}

% ************************************** %



% ^^^^^^^^^^^^^^^^^^^^^^^^^^^^^^^^^^^^^^ %
%               Question                 %
% ^^^^^^^^^^^^^^^^^^^^^^^^^^^^^^^^^^^^^^ %

\begin{enonce}
	En utilisant la relation \eqref{relation_v_V} en $x=0$, exprimer la constante en fonction de $q$ et $V_0$
\end{enonce}

\reponse{$qV_0$}

\begin{corrige}
	En $x=0$ la vitesse est nulle donc $
		\mathrm{C^{te}} = \frac{1}{2}mv^2(0) + qV(0) = qV \left( 1 - \left(\frac{0}{a}\right)^2\right) 
		= qV_0$.
\end{corrige}

% ************************************** %


\hauteurLargeurCadreReponse		{8mm}{4cm}


% ^^^^^^^^^^^^^^^^^^^^^^^^^^^^^^^^^^^^^^ %
%               Question                 %
% ^^^^^^^^^^^^^^^^^^^^^^^^^^^^^^^^^^^^^^ %

\begin{enonce}
	Exprimer $v(a)$ en fonction de $q$, $m$ et $V_0$
\end{enonce}
	
\reponse{$\sqrt{\frac{2qV_0}{m}}$}
	
\begin{corrige}
	On a
	\[
	qV_0 = \frac{1}{2}mv^2(a) + qV(a) = \frac{1}{2}mv^2(a) + qV\left( 1 - \left(\frac{a}{a}\right)^2\right) = \frac{1}{2}mv^2(a).
	\]
	Donc on a $v(a) = \sqrt{\frac{2qV_0}{m}}$.
\end{corrige}

% ************************************** %




% ^^^^^^^^^^^^^^^^^^^^^^^^^^^^^^^^^^^^^^ %
%               Question                 %
% ^^^^^^^^^^^^^^^^^^^^^^^^^^^^^^^^^^^^^^ %

\begin{enonce}
	Exprimer $v\left(\frac{a}{2}\right)$ en fonction de $q$, $m$ et $V_0$
\end{enonce}
	
\reponse{$\sqrt{\frac{qV_0}{2m}}$}
	
\begin{corrige}
On a 
	\[
		qV_0 = \frac{1}{2}mv^2\left(\frac{a}{2}\right) + qV_0\left(\frac{a}{2}\right) = \frac{1}{2}mv^2\left(\frac{a}{2}\right) + qV_0\left( 1 - \left(\frac{a}{2a}\right)^2\right) = \frac{1}{2}mv^2\left(\frac{a}{2}\right) +  \frac{3}{4}qV_0
	\]
	Donc, on a
	\[
		\frac{1}{2}mv^2\left(\frac{a}{2}\right) = \frac{1}{4}qV_0 \qquad \text{ et donc} \qquad	v\left(\frac{a}{2}\right) = \sqrt{\frac{qV_0}{2m}}.
	\]
\end{corrige}

% ************************************** %




% ^^^^^^^^^^^^^^^^^^^^^^^^^^^^^^^^^^^^^^ %
%               Question                 %
% ^^^^^^^^^^^^^^^^^^^^^^^^^^^^^^^^^^^^^^ %

\begin{enonce}
	Exprimer $v\left(\frac{a}{2}\right)$ en fonction de $v(a)$
\end{enonce}
	
\reponse{$\frac{v(a)}{2}$}
	
\begin{corrige}
On a 
	\[
		v\left(\frac{a}{2}\right) = \sqrt{\frac{qV_0}{2m}} = \sqrt{\frac{2qV_0}{4m}}  = \frac{1}{2}\sqrt{\frac{2qV_0}{m}} = \frac{v(a)}{2}.
	\]
\end{corrige}

% ************************************** %

% =============================================================================== %
\finEntrainement
% =============================================================================== %









% •••••••••••••••••••••••••••••••••••••••••••••••••••••••••••••••••••••••••••• %
% •••••••••••••••••••••••••••••••••••••••••••••••••••••••••••••••••••••••••••• %
\sectionFicheEntrainement{Du potentiel au champ électrique}
% •••••••••••••••••••••••••••••••••••••••••••••••••••••••••••••••••••••••••••• %
% •••••••••••••••••••••••••••••••••••••••••••••••••••••••••••••••••••••••••••• %



% ********************************************************************* %
%                           ENTRAÎNEMENT                                % 
% ********************************************************************* %
% ================ Métadonnées sur l'entraînement ===================== %

\titreEntrainementFacultatif	{Potentiel électrique dû à deux charges}
\hauteurLargeurCadreReponse		{7mm}{6cm}
\dureeResolutionFacultative		{4} % 1, 2, 3 ou 4
\typeEntrainement				{} % Newton ou QCM ou AN ou vide (exercice "normal")
\basiqueEtTransversal			{Y}
\nombreColonnesQuestions		{1} % vide, 1, 2, 3, etc.
\avecPlusieursQuestions			{Y} % Y ou N
\initialisationEntrainement
% ===================================================================== %

% mmmmmmmmmmmmmmmmmmmmmmmmmmmmmmmmmmmmmmmmmmmmmmmmmmmmmmmmm %
% 		  Insertion d'une image en regard d'un texte
% mmmmmmmmmmmmmmmmmmmmmmmmmmmmmmmmmmmmmmmmmmmmmmmmmmmmmmmmm %
% --------------------------------------------------------- %
\pourcentageDeLaPartieAGauche   {0.7}
% --------------------------------------------------------- %
%                                                           %
% -------------------- Partie à gauche -------------------- %
%                                                           %
                                \initialisationPartieGauche % 
%                                                           %
% 

Le potentiel électrique produit en un point $\ptM$ par une charge $q_1$ située en un point $\ptB$ est 
\[
V_1(\ptM) = \frac{1}{4\pi\varepsilon_0}\frac{q_1}{\ptB\ptM}.
\]

Afin d'obtenir les potentiels $V_1(\ptM)$ et $V_2(\ptM)$ créés par les charges $q_1$ et $q_2$ telles que $q = q_1 = -q_2$, ainsi que le potentiel total 
\[
V(\ptM) = V_1(\ptM) + V_2(\ptM),
\]
on cherche à exprimer les distances $\ptB\ptM$ et $\ptC\ptM$ en fonction des coordonnées $r$ et $\theta$ du point $\ptM$ et de la distance $a$ illustrées ci-contre.


% --------------------------------------------------------- %
%                                                           %
% -------------------- Partie à droite -------------------- %
%                                                           %
\initialisationPartieDroite %
%                                                           %
%                                                           %
\begin{center}
	\subimport{_images/}{fiche_MAG01_figure_potentiel_deux_charges_CBD_v2.tex}
\end{center}
% --------------------------------------------------------- %
                               \finalisationDuPartageDePage %					
% --------------------------------------------------------- %

\vspace{1mm}


Exprimer les grandeurs suivantes en fonction des paramètres indiqués.



% =============================================================================== %
\debutEntrainement
% =============================================================================== %


% ^^^^^^^^^^^^^^^^^^^^^^^^^^^^^^^^^^^^^^ %
%               Question                 %
% ^^^^^^^^^^^^^^^^^^^^^^^^^^^^^^^^^^^^^^ %

\begin{enonce}
	$\ptB\ptM$ en fonction de $x,y,a$. \minusMedskip
	
	{\itshape On pourra utiliser les coordonnées des points $\ptB$ et $\ptM$}
\end{enonce}

\reponse{$ \sqrt{ \left(x - a\right)^2 + y^2} $}

\begin{corrige}
	Dans le triangle $x\ptB\ptM$, on a  $\ptB\ptM =  \sqrt{ \left(x - a\right)^2 + y^2}$.
\end{corrige}

% ************************************** %



% ^^^^^^^^^^^^^^^^^^^^^^^^^^^^^^^^^^^^^^ %
%               Question                 %
% ^^^^^^^^^^^^^^^^^^^^^^^^^^^^^^^^^^^^^^ %

\begin{enonce}
\label{enonce_2}
	$r^2$ en fonction de $x,y$. \minusMedskip
	
	{\itshape On pourra chercher un triangle rectangle adéquat}
\end{enonce}

\reponse{$x^2+y^2$}

\begin{corrige}
	Dans le triangle $x\ptO\ptM$, on a   $r^2 = x^2 + y^2$.
\end{corrige}

% ************************************** %



% ^^^^^^^^^^^^^^^^^^^^^^^^^^^^^^^^^^^^^^ %
%               Question                 %
% ^^^^^^^^^^^^^^^^^^^^^^^^^^^^^^^^^^^^^^ %

\begin{enonce}
	$\ptB\ptM$ en fonction de $r,x,a$. \minusMedskip
	
	{\itshape On pourra utiliser  les réponses des questions a) et b)}
\end{enonce}

\reponse{$\sqrt{ r^2 - 2ax + a^2 }$}

\begin{corrige}
	En utilisant l'expression de $r^2$ en fonction de $x,y$, on a 
	\[
		\ptB\ptM =  \sqrt{ \left(x - a\right)^2 + y^2} =
		\sqrt{ x^2 + y^2 - 2ax + a^2} = \sqrt{ r^2 - 2ax + a^2 }.
	\]
\end{corrige}

% ************************************** %




% ^^^^^^^^^^^^^^^^^^^^^^^^^^^^^^^^^^^^^^ %
%               Question                 %
% ^^^^^^^^^^^^^^^^^^^^^^^^^^^^^^^^^^^^^^ %

\begin{enonce}
	$x$ en fontion de $r,\theta$.
\end{enonce}

\reponse{$r\cos(\theta)$}

\begin{corrige}
	Dans le triangle $x\ptO\ptM$, on a $\cos (\theta) = \frac{x}{r}$ et donc $x = r\cos(\theta)$.
\end{corrige}

% ************************************** %



% ^^^^^^^^^^^^^^^^^^^^^^^^^^^^^^^^^^^^^^ %
%               Question                 %
% ^^^^^^^^^^^^^^^^^^^^^^^^^^^^^^^^^^^^^^ %

\begin{enonce}
	$\ptB\ptM$ en fonction de $r,a,\theta$ \minusMedskip
	
	{\itshape On pourra utiliser  les réponses des questions c) et d)}
\end{enonce}

\reponse{$\sqrt{ r^2 - 2ar \cos (\theta) + a^2  }$}

\begin{corrige}
	En utilisant l'expression de $x$ en fonction de $r,\theta$, on peut écrire
	\[
		\ptB\ptM = \sqrt{r^2 - 2ax + a^2 }= \sqrt{ r^2  - 2ar \cos (\theta) + a^2 }.
	\]
\end{corrige}

% ************************************** %



% ^^^^^^^^^^^^^^^^^^^^^^^^^^^^^^^^^^^^^^ %
%               Question                 %
% ^^^^^^^^^^^^^^^^^^^^^^^^^^^^^^^^^^^^^^ %

\begin{enonce}
	$V_1$ en fonction de $q, r, a, \theta$.
\end{enonce}

\reponse{$\frac{1}{4\pi\varepsilon_0}\frac{q}{\sqrt{ r^2  - 2ar \cos (\theta) + a^2 }}$}

\begin{corrige}
	En utilisant les expressions de $V_1(\ptM)$ et de $\ptB\ptM$ en fonction de $r,a,\theta$,
	on a  
	\[
		V_1(\ptM) = \frac{1}{4\pi\varepsilon_0}\frac{q_1}{\ptB\ptM} = \frac{1}{4\pi\varepsilon_0}\frac{q}{\sqrt{ r^2  - 2ar \cos (\theta) + a^2 }}
	.\]
\end{corrige}

% ************************************** %




% ^^^^^^^^^^^^^^^^^^^^^^^^^^^^^^^^^^^^^^ %
%               Question                 %
% ^^^^^^^^^^^^^^^^^^^^^^^^^^^^^^^^^^^^^^ %

\begin{enonce}
	$\ptC\ptM$ en fonction de $x,y,a$. \minusMedskip
	
	{\itshape On pourra utiliser les coordonnées des points $\ptC$ et $\ptM$}
\end{enonce}

\reponse{$ \sqrt{ \left(x + a\right)^2 + y^2} $}

\begin{corrige}
	Dans le triangle $x\ptC\ptM$, on a $\ptC\ptM =  \sqrt{ \left(x + a\right)^2 + y^2}$.
\end{corrige}

% ************************************** %



% ^^^^^^^^^^^^^^^^^^^^^^^^^^^^^^^^^^^^^^ %
%               Question                 %
% ^^^^^^^^^^^^^^^^^^^^^^^^^^^^^^^^^^^^^^ %

\begin{enonce}
	$\ptC\ptM$ en fonction de  $r,x,a$. \minusMedskip
	
	{\itshape On pourra utiliser les réponses des questions b) et g)}
\end{enonce}

\reponse{$\sqrt{ r^2 + 2ax + a^2 }$}

\begin{corrige}
	En utilisant l'expression de $r^2$ en fonction  de $x,y$, on a 
	\[
		\ptC\ptM =  \sqrt{ \left(x + a\right)^2 + y^2} =
		\sqrt{ x^2 + y^2 + 2ax + a^2} = \sqrt{ r^2 + 2ax + a^2 }.
	\]
\end{corrige}

% ************************************** %



% ^^^^^^^^^^^^^^^^^^^^^^^^^^^^^^^^^^^^^^ %
%               Question                 %
% ^^^^^^^^^^^^^^^^^^^^^^^^^^^^^^^^^^^^^^ %

\begin{enonce}
	$\ptC\ptM$ en fonction de $r,a,\theta$. \minusMedskip
	
	{\itshape On pourra utiliser les réponses des questions d) et h)}
\end{enonce}

\reponse{$\sqrt{ r^2 + 2ar \cos (\theta) + a^2  }$}

\begin{corrige}
	En utilisant l'expression de $x$ en fonction de $r,\theta$, on a 
	$
		\ptC\ptM = \sqrt{r^2 + 2ax + a^2 }= \sqrt{ r^2  + 2ar \cos (\theta) + a^2 }
	$.
\end{corrige}

% ************************************** %




% ^^^^^^^^^^^^^^^^^^^^^^^^^^^^^^^^^^^^^^ %
%               Question                 %
% ^^^^^^^^^^^^^^^^^^^^^^^^^^^^^^^^^^^^^^ %

\begin{enonce}
	$V_2$ en fonction de $q, r, a, \theta$
\end{enonce}

\reponse{$-\frac{1}{4\pi\varepsilon_0}\frac{q}{\sqrt{ r^2  + 2ar \cos (\theta) + a^2 }}$}

\begin{corrige}
	En utilisant les expressions de $V_2(\ptM)$ et $\ptC\ptM$ en fonction de $r,a,\theta$, on a 
	\[
		V_2(\ptM) = \frac{1}{4\pi\varepsilon_0}\frac{q_2}{\ptC\ptM} = -\frac{1}{4\pi\varepsilon_0}\frac{q}{\sqrt{ r^2  + 2ar \cos (\theta) + a^2 }}.
		\]
\end{corrige}

% ************************************** %




% ^^^^^^^^^^^^^^^^^^^^^^^^^^^^^^^^^^^^^^ %
%               Question                 %
% ^^^^^^^^^^^^^^^^^^^^^^^^^^^^^^^^^^^^^^ %

\begin{enonce}
	$V$ en fonction de $q, r, a, \theta$
\end{enonce}

\reponse{\small
\begin{tabular}{l}
$\frac{1}{4\pi\varepsilon_0}q\Big(\frac{1}{\sqrt{ r^2  - 2ar \cos (\theta) + a^2 }}$ \\ 
\qquad\quad\quad${}-  \frac{1}{\sqrt{ r^2  + 2ar \cos (\theta) + a^2 }}  \Big)$
\end{tabular}}

\begin{corrige}
	En utilisant les expressions de $V_1(\ptM)$ et $V_2(\ptM)$, on trouve
	\[
		V(\ptM) = V_1(\ptM) + V_2(\ptM) = \frac{1}{4\pi\varepsilon_0}q\left(\frac{1}{\sqrt{ r^2  - 2ar \cos (\theta) + a^2 }} -  \frac{1}{\sqrt{ r^2  + 2ar \cos (\theta) + a^2 }}  \right).
	\]
\end{corrige}

% ************************************** %



% =============================================================================== %
\finEntrainement
% =============================================================================== %



\newpage


% ********************************************************************* %
%                           ENTRAÎNEMENT                                % 
% ********************************************************************* %
% ================ Métadonnées sur l'entraînement ===================== %

\titreEntrainementFacultatif	{Approximation de potentiels électriques}
\hauteurLargeurCadreReponse		{8mm}{5cm}
\dureeResolutionFacultative		{1} % 1, 2, 3 ou 4
\typeEntrainement				{} % Newton ou QCM ou AN ou vide (exercice "normal")
\basiqueEtTransversal			{Y}
\nombreColonnesQuestions		{1} % vide, 1, 2, 3, etc.
\avecPlusieursQuestions			{Y} % Y ou N
\initialisationEntrainement
% ===================================================================== %

Développer les expressions de potentiels électriques suivantes en calculant leur développement limité au voisinage de 0 à l'ordre indiqué et selon la variable spécifiée.

% =============================================================================== %
\debutEntrainement
% =============================================================================== %

% ^^^^^^^^^^^^^^^^^^^^^^^^^^^^^^^^^^^^^^ %
%               Question                 %
% ^^^^^^^^^^^^^^^^^^^^^^^^^^^^^^^^^^^^^^ %

\begin{enonce}
	À l'ordre 1 : $V\left(\frac{a}{r}\right) = \frac{1}{4\pi\varepsilon_0} \frac{q}{r} \left( 1 - \frac{a}{2r}\right)^{4}$
\end{enonce}

\reponse{$\frac{1}{4\pi\varepsilon_0} \frac{q}{r}\left( 1 - \frac{2a}{r}\right)$}

\begin{corrige}
	À l'ordre 1, on a  $\left(1+x\right)^{\alpha} \approx  1 + \alpha x$. Ainsi, on a  
	$
		V\left(\frac{a}{r}\right) \approx \frac{1}{4\pi\varepsilon_0} \frac{q}{r} \left( 1 - \frac{4a}{2r}\right) = \frac{1}{4\pi\varepsilon_0} \frac{q}{r}\left( 1 - \frac{2a}{r}\right)
	$.
\end{corrige}

% ************************************** %



% ^^^^^^^^^^^^^^^^^^^^^^^^^^^^^^^^^^^^^^ %
%               Question                 %
% ^^^^^^^^^^^^^^^^^^^^^^^^^^^^^^^^^^^^^^ %

\begin{enonce}
	À l'ordre 1 : $V\left(\frac{a}{r}\right) = \frac{1}{4\pi\varepsilon_0} \frac{q}{r} \left( \frac{1}{\sqrt{ 1 - \frac{a}{r} \cos (\theta)  }} - \frac{1}{\sqrt{ 1 + \frac{a}{r} \cos (\theta)  }} \right)$
\end{enonce}

\reponse{$\frac{1}{4\pi\varepsilon_0}\frac{qa\cos (\theta) }{r^2}$}

\begin{corrige}
	À l'ordre 1, on a $\left(1+x\right)^{\alpha} \approx  1 + \alpha x$. Ainsi, on a
	\[
		V\left(\frac{a}{r}\right) \approx \frac{1}{4\pi\varepsilon_0} \frac{q}{r} \left( 1  + \frac{a}{2r} \cos (\theta)  - \left( 1 -  \frac{a}{2r} \cos (\theta)  \right) \right) = \frac{1}{4\pi\varepsilon_0} \frac{qa\cos (\theta)}{r^2}.
	\]
\end{corrige}

% ************************************** %




% ^^^^^^^^^^^^^^^^^^^^^^^^^^^^^^^^^^^^^^ %
%               Question                 %
% ^^^^^^^^^^^^^^^^^^^^^^^^^^^^^^^^^^^^^^ %

\begin{enonce}
	À l'ordre 2 : $V\left(\theta\right) = \frac{1}{4\pi\varepsilon_0} \frac{qa\cos (\theta) }{r^2}$
\end{enonce}

\reponse{$\frac{1}{4\pi\varepsilon_0} \frac{qa}{r^2}\left(1 - \frac{1}{2}\theta^2\right)$}

\begin{corrige}
	À l'ordre 2, on a  $\cos(\theta) \approx  1 - \frac{1}{2}\theta^2$. Ainsi, on a 
	$
		V\left(\theta\right) \approx \frac{1}{4\pi\varepsilon_0} \frac{qa}{r^2}\left(1 - \frac{1}{2}\theta^2\right)
	$.
\end{corrige}

% ************************************** %




% ^^^^^^^^^^^^^^^^^^^^^^^^^^^^^^^^^^^^^^ %
%               Question                 %
% ^^^^^^^^^^^^^^^^^^^^^^^^^^^^^^^^^^^^^^ %

\begin{enonce}
	À l'ordre 1 : $V\left(\frac{a}{r}\right) = \frac{1}{4\pi\varepsilon_0} \frac{q}{r} \ln \left( 1 + \frac{a}{r} \right)$
\end{enonce}

\reponse{$\frac{1}{4\pi\varepsilon_0} \frac{qa}{r^2}$}

\begin{corrige}
	À l'ordre 1, on a $\ln\left(1+x\right) \approx  x$. Ainsi,  on a  
	$
		V\left(\frac{a}{r}\right) \approx \frac{1}{4\pi\varepsilon_0} \frac{q}{r} \frac{a}{r} = \frac{1}{4\pi\varepsilon_0} \frac{qa}{r^2}
	$.
\end{corrige}

% ************************************** %




% ^^^^^^^^^^^^^^^^^^^^^^^^^^^^^^^^^^^^^^ %
%               Question                 %
% ^^^^^^^^^^^^^^^^^^^^^^^^^^^^^^^^^^^^^^ %

\begin{enonce}
	À l'ordre 1 : $V\left(\frac{a}{r}\right) = \frac{1}{4\pi\varepsilon_0} \frac{q}{r} \ln \left( \frac{ \sqrt{1 + \frac{4a^2}{r^2}} + 1 }{ \sqrt{1 + \frac{4a^2}{r^2}} - 1 } \right)$
\end{enonce}

\reponse{$\frac{1}{4\pi\varepsilon_0} \frac{q}{r}\ln \left( 1 + \frac{r^2}{a^2} \right)$}

\begin{corrige}
	À l'ordre 1, on a $\left(1+x\right)^{\alpha} \approx  1 + \alpha x$. Ainsi, on a 
	\[
		V\left(\frac{a}{r}\right) \approx \frac{1}{4\pi\varepsilon_0} \frac{q}{r} \ln \left( \frac{ 1 + \frac{2a^2}{r^2} + 1 }{ 1 + \frac{2a^2}{r^2} - 1 } \right) = \frac{1}{4\pi\varepsilon_0} \frac{q}{r} \ln \left( \frac{2 + \frac{2a^2}{r^2}}{\frac{2a^2}{r^2}} \right) = \frac{1}{4\pi\varepsilon_0} \frac{q}{r}\ln \left( 1 + \frac{r^2}{a^2} \right).
	\]
\end{corrige}

% ************************************** %


% =============================================================================== %
\finEntrainement
% =============================================================================== %






% ********************************************************************* %
%                           ENTRAÎNEMENT                                % 
% ********************************************************************* %
% ================ Métadonnées sur l'entraînement ===================== %

\titreEntrainementFacultatif	{Calcul d'un champ électrique}
\hauteurLargeurCadreReponse		{7mm}{7cm}
\dureeResolutionFacultative		{2} % 1, 2, 3 ou 4
\typeEntrainement				{AN} % Newton ou QCM ou AN ou vide (exercice "normal")
\basiqueEtTransversal			{Y}
\nombreColonnesQuestions		{1} % vide, 1, 2, 3, etc.
\avecPlusieursQuestions			{Y} % Y ou N
\initialisationEntrainement
% ===================================================================== %

En coordonnées polaires, le champ $\vv{E}(\ptM)$ au point $\ptM$ s'exprime en fonction du potentiel $V(\ptM)$ par la formule
	\[
		\vv{E}(\ptM) = -\diffp{V(\ptM)}{r}\vv{e_{r}} -\frac{1}{r}\diffp{V(\ptM)}{\theta}\vv{e_{\theta}}.
	\]
On donne 
\[
\varepsilon_0 = \np[C.V^{-1}.m^{-1}]{8,85e-12}, \quad q=\np[C]{6,0e-11}\quad\text{et}\quad a = \np[mm]{4,0}.
\]

\medskip

Dans cet entraînement, on suppose que $V(\ptM) = \frac{1}{4\pi\varepsilon_0} \frac{q \sin\left( 2\theta\right) }{r}$.


% =============================================================================== %
\debutEntrainement
% =============================================================================== %

% ^^^^^^^^^^^^^^^^^^^^^^^^^^^^^^^^^^^^^^ %
%               Question                 %
% ^^^^^^^^^^^^^^^^^^^^^^^^^^^^^^^^^^^^^^ %

\begin{enonce}
	Exprimer $\vv{E}(\ptM)$ 
\end{enonce}

\reponse{$\frac{1}{4\pi\varepsilon_0} \frac{q}{r^2} \left( \sin\left( 2\theta\right) \vv{e_{r}} -  2\cos\left( 2\theta\right) \vv{e_{\theta}} \right)$}

\begin{corrige}
On calcule :
	\begin{align*}
		\vv{E}(\ptM) &= -\diffp{}{r} \left(  \frac{1}{4\pi\varepsilon_0} \frac{q \sin\left( 2\theta\right) }{r} \right)  \vv{e_{r}} - \frac{1}{r}\diffp{}{\theta} \left(  \frac{1}{4\pi\varepsilon_0} \frac{q \sin\left( 2\theta\right) }{r} \right)  \vv{e_{\theta}}\\
		%\vv{E}(\ptM) 
		&= -\frac{1}{4\pi\varepsilon_0} q \left( \sin\left( 2\theta\right) \diffp{}{r}\left( \frac{1}{r}\right)\vv{e_{r}} + \frac{1}{r^2}\diffp{\sin\left( 2\theta\right)}{\theta}\vv{e_{\theta}}  \right)\\
		%\vv{E}(\ptM) 
		&= \frac{1}{4\pi\varepsilon_0}\frac{q}{r^2} \left( \sin\left( 2\theta\right) \vv{e_{r}} - 2\cos\left( 2\theta\right) \vv{e_{\theta}} \right).
	\end{align*}
\end{corrige}

% ************************************** %



% ^^^^^^^^^^^^^^^^^^^^^^^^^^^^^^^^^^^^^^ %
%               Question                 %
% ^^^^^^^^^^^^^^^^^^^^^^^^^^^^^^^^^^^^^^ %

\begin{enonce}
	Exprimer $\vv{E}(\ptM)$ pour $\ptM\left(r=\frac{a}{2},\theta=\pi\right)$.
\end{enonce}

\reponse{$-\frac{8}{4\pi\varepsilon_0} \frac{q}{a^2} \vv{e_{\theta}}$}

\begin{corrige}
	Pour $\ptM\left(r=\frac{a}{2},\theta=\pi\right)$ le champ est
	$
		\vv{E}(\ptM) = \frac{1}{4\pi\varepsilon_0} \frac{q}{\left(\frac{a}{2}\right)^2} \left( \sin\left( 2\pi\right) \vv{e_{r}} -  2\cos\left( 2\pi\right) \vv{e_{\theta}} \right) = -\frac{8}{4\pi\varepsilon_0} \frac{q}{a^2} \vv{e_{\theta}}
	$.
\end{corrige}

% ************************************** %



% ^^^^^^^^^^^^^^^^^^^^^^^^^^^^^^^^^^^^^^ %
%               Question                 %
% ^^^^^^^^^^^^^^^^^^^^^^^^^^^^^^^^^^^^^^ %

\begin{enonce}
	À l'aide des données, calculer $\| \vv{E}(\ptM) \|$ en V.m$^{-1}$.
\end{enonce}

\reponse{$\np[V.m^{-1}]{2,7e5}$}

\begin{corrige}
On a
	\[
		\| \vv{E}(\ptM) \| = \frac{8}{4\pi\varepsilon_0} \frac{q}{a^2} = \frac{8}{ 4\pi \times \np[C.V^{-1}.m^{-1}]{8,85e-12} } \frac{\np[C]{6,0e-11}}{\left(\np[m]{4,0e-3}\right)^2} = \np[V.m^{-1}]{2,7e5}.
	\]
\end{corrige}

% ************************************** %

% =============================================================================== %
\finEntrainement
% =============================================================================== %






% ********************************************************************* %
%                           ENTRAÎNEMENT                                % 
% ********************************************************************* %
% ================ Métadonnées sur l'entraînement ===================== %

\titreEntrainementFacultatif	{\textit{Bis repetita}}
\hauteurLargeurCadreReponse		{7mm}{7cm}
\dureeResolutionFacultative		{2} % 1, 2, 3 ou 4
\typeEntrainement				{AN} % Newton ou QCM ou AN ou vide (exercice "normal")
\basiqueEtTransversal			{Y}
\nombreColonnesQuestions		{1} % vide, 1, 2, 3, etc.
\avecPlusieursQuestions			{Y} % Y ou N
\initialisationEntrainement
% ===================================================================== %

On reprend l'entraînement précédent avec les mêmes données, mais un potentiel électrique différent. 

Dans cet entraînement, on suppose que $V(\ptM) = \frac{1}{4\pi\varepsilon_0} \frac{qa\cos (\theta)}{r^2}$.

% =============================================================================== %
\debutEntrainement
% =============================================================================== %

% ^^^^^^^^^^^^^^^^^^^^^^^^^^^^^^^^^^^^^^ %
%               Question                 %
% ^^^^^^^^^^^^^^^^^^^^^^^^^^^^^^^^^^^^^^ %
\begin{enonce}
	Exprimer $\vv{E}(\ptM)$
\end{enonce}

\reponse{$\frac{1}{4\pi\varepsilon_0}\frac{qa}{r^3} \left( 2\cos(\theta) \vv{e_{r}} + \sin (\theta) \vv{e_{\theta}} \right)$}

\begin{corrige}
On calcule :
	\begin{align*}
		\vv{E}(\ptM) &= -\diffp{}{r} \left(  \frac{1}{4\pi\varepsilon_0}\frac{qa\cos (\theta) }{r^2} \right)  \vv{e_{r}} - \frac{1}{r}\diffp{}{\theta} \left(  \frac{1}{4\pi\varepsilon_0}\frac{qa\cos (\theta) }{r^2} \right)  \vv{e_{\theta}}\\
		%\vv{E}(\ptM) 
		&= -\frac{1}{4\pi\varepsilon_0}qa \left( \cos (\theta) \diffp{}{r}\left( \frac{1}{r^2}\right)\vv{e_{r}} + \frac{1}{r^3}\diffp{\cos (\theta)}{\theta}\vv{e_{\theta}}  \right)\\
		%\vv{E}(\ptM) 
		&= \frac{1}{4\pi\varepsilon_0}\frac{qa}{r^3} \left( 2\cos(\theta) \vv{e_{r}} + \sin (\theta) \vv{e_{\theta}} \right).
	\end{align*}
\end{corrige}
% ************************************** %

% ^^^^^^^^^^^^^^^^^^^^^^^^^^^^^^^^^^^^^^ %
%               Question                 %
% ^^^^^^^^^^^^^^^^^^^^^^^^^^^^^^^^^^^^^^ %
\begin{enonce}
	Exprimer $\vv{E}(\ptM)$ pour $\ptM\left(r=a,\theta=\frac{\pi}{2}\right)$.
\end{enonce}

\reponse{$\frac{1}{4\pi\varepsilon_0} \frac{q}{a^2} \vv{e_{\theta}}$}

\begin{corrige}
	Pour $\ptM\left(r=a,\theta=\frac{\pi}{2}\right)$ le champ est
	\[
		\vv{E}(\ptM) = \frac{1}{4\pi\varepsilon_0} \frac{qa}{a^3} \left( 2\cos\left(\frac{\pi}{2}\right) \vv{e_{r}} +  \sin\left( \frac{\pi}{2}\right) \vv{e_{\theta}} \right) = \frac{1}{4\pi\varepsilon_0} \frac{q}{a^2} \vv{e_{\theta}}.
	\]
\end{corrige}
% ************************************** %

% ^^^^^^^^^^^^^^^^^^^^^^^^^^^^^^^^^^^^^^ %
%               Question                 %
% ^^^^^^^^^^^^^^^^^^^^^^^^^^^^^^^^^^^^^^ %
\begin{enonce}
	À l'aide des données, calculer $\| \vv{E}(\ptM) \|$ en V.m$^{-1}$.
\end{enonce}

\reponse{$\np[V.m^{-1}]{3,4e4}$}

\begin{corrige}
On a 
	\[
		\| \vv{E}(\ptM) \| = \frac{1}{4\pi\varepsilon_0} \frac{q}{a^2} = \frac{1}{ 4\pi \times \np[C.V^{-1}.m^{-1}]{8,85e-12} } \frac{\np[C]{6,0e-11}}{\left(\np[m]{4,0e-3}\right)^2} = \np[V.m^{-1}]{3,4e4}.
	\]
\end{corrige}
% ************************************** %

% =============================================================================== %
\finEntrainement
% =============================================================================== %





\newpage



% •••••••••••••••••••••••••••••••••••••••••••••••••••••••••••••••••••••••••••• %
% •••••••••••••••••••••••••••••••••••••••••••••••••••••••••••••••••••••••••••• %
\sectionFicheEntrainement{Du champ au potentiel électrique}
% •••••••••••••••••••••••••••••••••••••••••••••••••••••••••••••••••••••••••••• %
% •••••••••••••••••••••••••••••••••••••••••••••••••••••••••••••••••••••••••••• %



% ********************************************************************* %
%                           ENTRAÎNEMENT                                % 
% ********************************************************************* %
% ================ Métadonnées sur l'entraînement ===================== %

\titreEntrainementFacultatif	{Champ électrique produit par un condensateur}
\hauteurLargeurCadreReponse		{7mm}{7cm}
\dureeResolutionFacultative		{3} % 1, 2, 3 ou 4
\typeEntrainement				{} % Newton ou QCM ou AN ou vide (exercice "normal")
\basiqueEtTransversal			{Y}
\nombreColonnesQuestions		{1} % vide, 1, 2, 3, etc.
\avecPlusieursQuestions			{Y} % Y ou N
\initialisationEntrainement
% ===================================================================== %

Un condensateur produit un champ $\vv{E} = E(x)\vv{e_x}$ entre ses deux armatures positionnées en $x = 0$ et $x = d$. La différence de potentiels entre les armatures est liée au champ de telle manière que \vspace{-1mm}
\[
	V(0) - V(d) = \int_{0}^{d} E(x) \d{}x.\vspace{-1mm}
\]
On considère que l'armature en $x=d$ est la masse du circuit, son potentiel est donc considéré comme nul. 

\medskip

Exprimer le potentiel $V(0)$ pour les différentes formes de champ $E(x)$.

% =============================================================================== %
\debutEntrainement
% =============================================================================== %

% ^^^^^^^^^^^^^^^^^^^^^^^^^^^^^^^^^^^^^^ %
%               Question                 %
% ^^^^^^^^^^^^^^^^^^^^^^^^^^^^^^^^^^^^^^ %

\begin{enonce}
	$E(x) = E_0 \left( 1 - \frac{x}{d} \right)$
\end{enonce}

\reponse{$\frac{1}{2}E_0d$}

\begin{corrige}
On a $V(0) = \int_{0}^{d} E_0 \left( 1 - \frac{x}{d} \right) \d{}x = E_0 \left[ -\frac{d}{2}\left( 1 - \frac{x}{d} \right)^2  \right]^{d}_{0} = \frac{1}{2}E_0d$.
\end{corrige}

% ************************************** %


% ^^^^^^^^^^^^^^^^^^^^^^^^^^^^^^^^^^^^^^ %
%               Question                 %
% ^^^^^^^^^^^^^^^^^^^^^^^^^^^^^^^^^^^^^^ %

\begin{enonce}
	$E(x) = E_0 \left( 1 - \frac{x}{d} \right)^2$
\end{enonce}

\reponse{$\frac{1}{3}E_0d$}

\begin{corrige}
On a $V(0) = \int_{0}^{d} E_0 \left( 1 - \frac{x}{d} \right)^2 \d{}x = E_0 \left[ -\frac{d}{3}\left( 1 - \frac{x}{d} \right)^3  \right]^{d}_{0} = \frac{1}{3}E_0d$.
\end{corrige}

% ************************************** %


% ^^^^^^^^^^^^^^^^^^^^^^^^^^^^^^^^^^^^^^ %
%               Question                 %
% ^^^^^^^^^^^^^^^^^^^^^^^^^^^^^^^^^^^^^^ %

\begin{enonce}
	$E(x) = E_0 \sin\left( \frac{3\pi}{2}\frac{x}{d} \right)$
\end{enonce}

\reponse{$\frac{2}{3\pi}E_0d$}

\begin{corrige}
On a $V(0) = \int_{0}^{d} E_0 \sin\left( \frac{3\pi}{2}\frac{x}{d} \right) \d{}x = E_0 \left[ -\frac{2d}{3\pi}\cos\left( \frac{3\pi}{2}\frac{x}{d} \right)  \right]^{d}_{0} = \frac{2}{3\pi}E_0d$.
\end{corrige}

% ************************************** %


% ^^^^^^^^^^^^^^^^^^^^^^^^^^^^^^^^^^^^^^ %
%               Question                 %
% ^^^^^^^^^^^^^^^^^^^^^^^^^^^^^^^^^^^^^^ %
\begin{enonce}
	$E(x) = E_0 \left( 1 - \eEuler^{-x/d} \right)$
\end{enonce}

\reponse{$E_0de^{-1}$}


%E_0 d\left(2-\eEuler^{-1}\right)

\begin{corrige}
On a $V(0) = \int_{0}^{d} E_0 \left( 1 - \eEuler^{-x/d} \right) \d{}x =
		E_0 \Big[ x  \Big]^{d}_{0} - E_0 \Big[ -d\eEuler^{-x/d}  \Big]^{d}_{0} = E_0d\eEuler^{-1}$.
\end{corrige}

% ************************************** %

% =============================================================================== %
\finEntrainement
% =============================================================================== %






% •••••••••••••••••••••••••••••••••••••••••••••••••••••••••••••••••••••••••••• %
% •••••••••••••••••••••••••••••••••••••••••••••••••••••••••••••••••••••••••••• %
\sectionFicheEntrainement{Distributions continues de charges}
% •••••••••••••••••••••••••••••••••••••••••••••••••••••••••••••••••••••••••••• %
% •••••••••••••••••••••••••••••••••••••••••••••••••••••••••••••••••••••••••••• %





% ********************************************************************* %
%                           ENTRAÎNEMENT                                % 
% ********************************************************************* %
% ================ Métadonnées sur l'entraînement ===================== %

\titreEntrainementFacultatif	{Charge d'une sphère}
\hauteurLargeurCadreReponse		{7mm}{8cm}
\dureeResolutionFacultative		{4} % 1, 2, 3 ou 4
\typeEntrainement				{} % Newton ou QCM ou AN ou vide (exercice "normal")
\basiqueEtTransversal			{Y}
\nombreColonnesQuestions		{1} % vide, 1, 2, 3, etc.
\avecPlusieursQuestions			{Y} % Y ou N
\initialisationEntrainement
% ===================================================================== %


On souhaite déterminer la charge électrique totale $Q$ contenue dans une sphère de rayon $R$ et de densité de charges $\rho\left(r,\theta,\varphi\right)$. Pour ce faire on doit intégrer la densité de charges sur toute la sphère $\mathcal{S}$ en utilisant la formule \vspace{-1mm}
\[
Q = \iiint_{\mathcal{S}} \rho\left(r,\theta,\varphi\right) \ \d{}\tau.
\]

On peut démontrer que 
le volume d'intégration élémentaire est $\d{}\tau=r^2 \sin (\theta) \, \d{}r \, \d{} \theta \, \d{} \varphi$ avec pour une sphère $r \in \left[ 0,\, R\right]$, $\theta \in \left[ 0,\, \pi\right]$ et $\varphi \in \left[ 0,\, 2\pi\right]$. Ainsi, on a 
\[
Q = \int_{0}^{2\pi} \int_{0}^{\pi} \int_{0}^{R} \rho\left(r,\theta,\varphi\right) r^2 \sin (\theta) \, \d{}r \, \d{} \theta \, \d{} \varphi.
\]



Exprimer la charge électrique totale $Q$ contenue dans la sphère en fonction de son rayon $R$ pour les différentes densités de charges suivantes.


% =============================================================================== %
\debutEntrainement
% =============================================================================== %

% ^^^^^^^^^^^^^^^^^^^^^^^^^^^^^^^^^^^^^^ %
%               Question                 %
% ^^^^^^^^^^^^^^^^^^^^^^^^^^^^^^^^^^^^^^ %

\begin{enonce}
	$\rho(r,\theta,\varphi) = 2 \rho_0$ 
\end{enonce}

\reponse{$\frac{8}{3}\pi R^3 \rho_0$}

\begin{corrige}
On calcule
\begin{align*}
	Q &= \int_{0}^{2\pi} \int_{0}^{\pi} \int_{0}^{R} 2\rho_0  r^2 \sin (\theta) 
	 \d{}r \d{} \theta \d{} \varphi \\
	 &= 2\rho_0 \int_{0}^{2\pi} \d{} \varphi \int_{0}^{\pi} \sin (\theta) \d{} \theta \int_{0}^{R} r^2 \d{}r
	 = 2\rho_0 \Big[ \varphi \Big]^{2\pi}_{0} \times \Big[ -\cos (\theta) \Big]^{\pi}_{0} \times \Big[ \frac{r^3}{3} \Big]^{R}_{0} \\
	 &= 2\rho_0 \left( 2\pi - 0 \right) \left( -\cos \pi + \cos 0 \right) \left( \frac{R^3}{3} - \frac{0}{3} \right)
	 = 2\rho_0 \times 2 \pi \times 2 \times \frac{R^3}{3} = \frac{8}{3}\pi R^3 \rho_0.
\end{align*}
\end{corrige}

% ************************************** %



% ^^^^^^^^^^^^^^^^^^^^^^^^^^^^^^^^^^^^^^ %
%               Question                 %
% ^^^^^^^^^^^^^^^^^^^^^^^^^^^^^^^^^^^^^^ %

\begin{enonce}	
	$\rho(r,\theta,\varphi) = 2\left(\frac{r}{R}\right)^2 \rho_0$
\end{enonce}

\reponse{$\frac{8}{5}\pi R^3 \rho_0$}

\begin{corrige}
On calcule
	\begin{align*}
		Q &= 2 \int_{0}^{2\pi} \d{} \varphi \int_{0}^{\pi} \sin (\theta) \d{} \theta \int_{0}^{R} \left( \frac{r}{R} \right)^2\rho_0 r^2 \d{}r
		 = 2\rho_0 \times 2 \pi \times 2 \times \int_{0}^{R} \frac{r^4}{R^2} \d{}r \\
		 &= 8\rho_0\pi \left[ \frac{1}{5} \frac{r^5}{R^2} \right]^{R}_{0} = 8\pi \left( \frac{1}{5} \frac{R^5}{R^2} - \frac{1}{5} \frac{0}{R^2} \right)\rho_0 = \frac{8}{5}\pi R^3\rho_0.
	\end{align*}
\end{corrige}

% ************************************** %



% ^^^^^^^^^^^^^^^^^^^^^^^^^^^^^^^^^^^^^^ %
%               Question                 %
% ^^^^^^^^^^^^^^^^^^^^^^^^^^^^^^^^^^^^^^ %
\begin{enonce}	
	$\rho(r,\theta,\varphi) = 2\left(\frac{r}{R}\right)^2\sin\left(\frac{\varphi}{2}\right)\rho_0$
\end{enonce}

\reponse{$\frac{16}{5}R^3\rho_0$}

\begin{corrige}
On calcule
	\begin{align*}
		Q &= 2 \int_{0}^{2\pi}  \sin \left( \frac{\varphi}{2} \right)\rho_0 \d{} \varphi \int_{0}^{\pi} \sin (\theta) \d{} \theta \int_{0}^{R} \left( \frac{r}{R} \right)^2 r^2 \d{}r
		 = 2\rho_0 \times \left[ -2 \cos \frac{\varphi}{2} \right]^{2\pi}_{0} \times 2 \times \frac{1}{5}R^3\\
		 &= \frac{4}{5}R^3 \left( -2 \cos \pi + 2 \cos 0  \right)\rho_0
		 = \frac{4}{5}R^3 \left( 2+ 2 \right)\rho_0 = \frac{16}{5}R^3\rho_0.
	\end{align*}
\end{corrige}
% ************************************** %

% =============================================================================== %
\finEntrainement
% =============================================================================== %



\newpage


% ********************************************************************* %
%                           ENTRAÎNEMENT                                % 
% ********************************************************************* %
% ================ Métadonnées sur l'entraînement ===================== %

\titreEntrainementFacultatif	{Charge d'un cylindre}
\hauteurLargeurCadreReponse		{7mm}{8cm}
\dureeResolutionFacultative		{4} % 1, 2, 3 ou 4
\typeEntrainement				{} % Newton ou QCM ou AN ou vide (exercice "normal")
\basiqueEtTransversal			{Y}
\nombreColonnesQuestions		{1} % vide, 1, 2, 3, etc.
\avecPlusieursQuestions			{Y} % Y ou N
\initialisationEntrainement
% ===================================================================== %

% mmmmmmmmmmmmmmmmmmmmmmmmmmmmmmmmmmmmmmmmmmmmmmmmmmmmmmmmm %
% 		  Insertion d'une image en regard d'un texte
% mmmmmmmmmmmmmmmmmmmmmmmmmmmmmmmmmmmmmmmmmmmmmmmmmmmmmmmmm %
% --------------------------------------------------------- %
\pourcentageDeLaPartieAGauche   {0.55}
% --------------------------------------------------------- %
%                                                           %
% -------------------- Partie à gauche -------------------- %
%                                                           %
                                \initialisationPartieGauche % 
%                                                           %
%                                                           %
On souhaite déterminer la charge électrique totale $Q$ contenue dans un cylindre de rayon $R$, de hauteur $h$ et de densité de charges $\rho\left(r,\theta,z\right)$. 
Pour ce faire on doit intégrer la densité de charges sur tout le cylindre $\mathcal{C}$. 

Comme on peut le voir sur la figure ci-contre, le volume d'intégration élémentaire est 
\[
\d{}\tau=r \, \d{}r \, \d{} \theta \, \d{} z
\]
avec pour un cylindre $r \in \left[ 0,\,R\right]$, $\theta \in \left[ 0,\,2\pi\right]$ et $z \in \left[ 0,\,h\right]$. 

Ainsi, on a \vspace{-1.5mm}
\[
Q = \int_{0}^{h} \int_{0}^{2\pi} \int_{0}^{R} \rho\left(r,\theta,z\right) r \, \d{}r \, \d{} \theta \, \d{} z.
\]


% --------------------------------------------------------- %
%                                                           %
% -------------------- Partie à droite -------------------- %
%                                                           %
                                \initialisationPartieDroite %
%                                                           %
%                                                           %
\begin{center}
	\subimport{_images/}{fiche_MAG01_figure_integration_cylindrique_CBD_v2.tex}
\end{center}
% --------------------------------------------------------- %
                               \finalisationDuPartageDePage %					
% --------------------------------------------------------- %


\medskip

Exprimer la charge électrique totale $Q$ contenue dans le cylindre en fonction de son rayon $R$ et de sa hauteur $h$ pour les différentes densités de charges suivantes.


% =============================================================================== %
\debutEntrainement
% =============================================================================== %




% ^^^^^^^^^^^^^^^^^^^^^^^^^^^^^^^^^^^^^^ %
%               Question                 %
% ^^^^^^^^^^^^^^^^^^^^^^^^^^^^^^^^^^^^^^ %
\begin{enonce}
	$\rho\left(r,\theta,z\right) = 3$ 
\end{enonce}

\reponse{$3\pi R^2h$}


\begin{corrige}
On calcule
\begin{align*}
	Q &= \int_{0}^{h} \int_{0}^{2\pi} \int_{0}^{R} 3 r 
	\d{}r \d{} \theta \d{} z
	= 3 \int_{0}^{h} \d{} z \int_{0}^{2\pi} \d{} \theta \int_{0}^{R} r \d{}r
	= 3 \Big[ z \Big]^{h}_{0}\times \Big[ \theta \Big]^{2\pi}_{0} \times \Big[ \frac{r^2}{2} \Big]^{R}_{0}\\
	&= 3 \left( h -0 \right) \left( 2\pi - 0 \right) \left( \frac{R^2}{2} - 0 \right)
	= 3\pi R^2h.
\end{align*}
\end{corrige}
% ************************************** %

% ^^^^^^^^^^^^^^^^^^^^^^^^^^^^^^^^^^^^^^ %
%               Question                 %
% ^^^^^^^^^^^^^^^^^^^^^^^^^^^^^^^^^^^^^^ %
\begin{enonce}	
	$\rho\left(r,\theta,z\right) = 2\left(\frac{r}{R}\right)^3$
\end{enonce}

\reponse{$\frac{4}{5}\pi R^2h$}

\begin{corrige}
On calcule
	\begin{align*}
		Q &= 2\int_{0}^{h} \d{} z \int_{0}^{2\pi} \d{} \theta  \int_{0}^{R} \left( \frac{r}{R} \right)^3 r \d{}r 
		= 2 \times h \times 2\pi \int^{R}_{0} \frac{r^4}{R^3} \d{}r \\
		&= 4\pi h \left[ \frac{1}{5} \frac{r^5}{R^3} \right]^{R}_{0}
		= 4\pi h \left( \frac{1}{5} \frac{R^5}{R^3} - \frac{1}{5} \frac{0}{R^3} \right)
		= \frac{4}{5}\pi R^2h.
	\end{align*}
\end{corrige}
% ************************************** %

% ^^^^^^^^^^^^^^^^^^^^^^^^^^^^^^^^^^^^^^ %
%               Question                 %
% ^^^^^^^^^^^^^^^^^^^^^^^^^^^^^^^^^^^^^^ %
\begin{enonce}	
	$\rho\left(r,\theta,z\right) = 2\left(\frac{r}{R}\right)^3\left(\frac{z}{h}\right)^2\sin\left(\frac{\theta}{2}\right)$
\end{enonce}

\reponse{$\frac{8}{15}R^2h$}

\begin{corrige}
On calcule
	\begin{align*}
		Q &= 2\int_{0}^{h} \left( \frac{z}{h} \right)^{2} \d{} z \int_{0}^{2\pi} \sin \left( \frac{\theta}{2} \right) \d{} \theta  \int_{0}^{R} \left( \frac{r}{R} \right)^3 r \d{}r 
		= 2 \times \frac{1}{5}R^2 \left[ \frac{1}{3}\frac{z^3}{h^2} \right]^{h}_{0} \left[ -2 \cos \frac{\theta}{2} \right]^{2\pi}_{0}\\
		&= \frac{2}{5}R^2 \left( \frac{1}{3}\frac{h^3}{h^2} - \frac{1}{3}\frac{0}{h^2} \right) \left( -2\cos \pi + 2 \cos 0 \right)
		= \frac{2}{5}R^2 \times \frac{1}{3} h \times 4 = \frac{8}{15}R^2h.
	\end{align*}
\end{corrige}
% ************************************** %

% =============================================================================== %
\finEntrainement
% =============================================================================== %





% ---------------------------------------------------------------------------- %
%                       Affichage des réponses mélangées                       %
% ---------------------------------------------------------------------------- %
\afficheReponsesMelangees
% ---------------------------------------------------------------------------- %

%%%%%%%%%%%%%%%%%%%%%%%%%%%%%%%%%%%%%%%%%%%%%%%%%%%%%%%%%%%%%%%%%%%%%%%%%%%%%%%%%%%%%%%%%%%%%%
\finFicheEntrainement                                                                        %
%%%%%%%%%%%%%%%%%%%%%%%%%%%%%%%%%%%%%%%%%%%%%%%%%%%%%%%%%%%%%%%%%%%%%%%%%%%%%%%%%%%%%%%%%%%%%%


% ======================================================= % 
% ============== Gestion de la compilation ============== %
% ======================================================= % 
\ifdefined\mainIsLoaded\else
\printReponsesEtCorriges
\end{document}\fi
% ======================================================= % 
% ======================================================= % 